% Vietnamese translation for OI Public Library
% Source-File: IOI中国国家候选队论文/2005/任恺/任恺.ppt
\documentclass[11pt]{article}
\usepackage{oipl-vn}

\title{Bản trình chiếu: Tư tưởng cơ bản của lý thuyết đồ thị}
\author{Ren Kai}
\date{2005}

\begin{document}
\maketitle

\origtitle{Slide companion for graph-theory methods}
\sourcepath{IOI China candidate team papers / 2005 / Ren Kai / Ren Kai.ppt}

\section{Phạm vi bản dịch}

Bản trình chiếu đi kèm bài viết về các tư tưởng và phương pháp cơ bản của lý
thuyết đồ thị. Phần chữ khôi phục được từ PPT cũ có nhắc tới POI 2002, các
ký hiệu cạnh \(E\), đỉnh \(u,v\), và độ phức tạp \(O(nm)\). Nội dung dưới đây
dựa trên bản dịch bài viết cùng thư mục để ghi lại mạch trình bày.

\section{Mạch nội dung}

Tư tưởng chính là biến quan hệ rời rạc trong đề bài thành cấu trúc đồ thị:
đỉnh biểu diễn đối tượng, cạnh biểu diễn ràng buộc, còn đường đi hoặc thành
phần liên thông biểu diễn khả năng lan truyền thông tin. Khi mô hình đúng,
nhiều bài toán chuyển về tìm cầu, khớp, thành phần liên thông mạnh, hoặc cây
khung.

\section{Thông điệp thuật toán}

Các slide nhấn mạnh rằng không nên chỉ nhớ tên thuật toán. Cần nhận ra vì sao
đồ thị mô tả đúng quan hệ trong dữ liệu, sau đó mới chọn DFS, BFS, Tarjan,
Dijkstra hay mô hình luồng phù hợp. Bản trình chiếu dùng các ví dụ thi để
minh họa cách giảm bài toán gốc về một bất biến hoặc một cấu trúc đồ thị nhỏ
hơn.

\end{document}
